\documentclass[12pt]{article}
\usepackage[utf8]{inputenc}
\usepackage[T1]{fontenc}
\usepackage{hyperref}
\title{Tlacoapa Me'phaa Language Resources for K-12 Curriculum Development}
\author{Compiled by Claude Sonnet 5 (claude-sonnet-5) for Andr\'as Kornai}
\date{}
\begin{document}
\maketitle
\section{Introduction}
Tlacoapa Me'phaa (ISO 639-3 tpl, Glottocode tlac1237) is an Otomanguean language of Guerrero, Mexico. This report is based on a single round of web searches. Little material specific to the Tlacoapa variety was found, and no school curriculum, primer or textbook for it was found.
\section{Primary Educational Resources}
No primers, textbooks or school programmes specific to Tlacoapa Me'phaa were found in the search results.
\section{Linguistic and Reference Materials}
A University of Kansas news item describes a talking-dictionary project for Me'phaa (Guerrero), with Me'phaa--English--Spanish content. It was not confirmed to cover the Tlacoapa variety specifically: \url{https://news.ku.edu/news/article/2021/04/29/talking-dictionary-link-lawrence-kansas-and-guerrero-mexico}

Glottolog language page for the Me'phaa group (classification only): \url{https://glottolog.org/resource/languoid/id/meph1234}

Wikitongues language page: \url{https://wikitongues.org/languages/tpl/}

The Endangered Languages Alliance page on Tlapanec mentions general lexical and grammatical resources for Tlapanec varieties; it is not specific to Tlacoapa: \url{https://www.elalliance.org/languages/tlapanec}
\section{Religious and Community Resources}
ScriptureEarth lists Scripture resources for tpl (New Testament audio and related media): \url{https://www.scriptureearth.org/00eng.php?iso=tpl}
\section{Digital and Audio Resources}
The ScriptureEarth page above is the only audio source found. The Endangered Languages Alliance resource list may be worth checking: \url{https://endangeredlanguages.com/hi/node/38942/resources}
\section{Curriculum Development Context}
The search results show only scripture, a dictionary project and general reference pages. A curriculum would have to start from community work and from the linguistic documentation that exists, since no existing classroom materials were found.
\end{document}